\documentclass[11pt]{article}

\usepackage[a4paper,margin=1in]{geometry}
\usepackage{amsmath,amssymb,amsthm}
\usepackage{mathtools}
\usepackage{xcolor}
\usepackage{hyperref}
\usepackage{cleveref}
\usepackage{microtype}
\usepackage{enumitem}

\hypersetup{
  colorlinks=true,
  linkcolor=blue!50!black,
  citecolor=blue!50!black,
  urlcolor=blue!50!black
}

\theoremstyle{plain}
\newtheorem{theorem}{Theorem}[section]
\newtheorem{proposition}[theorem]{Proposition}
\newtheorem{lemma}[theorem]{Lemma}
\newtheorem{corollary}[theorem]{Corollary}

\theoremstyle{definition}
\newtheorem{definition}[theorem]{Definition}
\newtheorem{example}[theorem]{Example}

\theoremstyle{remark}
\newtheorem{remark}[theorem]{Remark}

\newcommand{\HH}{\mathcal{H}}
\newcommand{\HtH}{H^{2}}
\newcommand{\KB}{K_{B}}
\newcommand{\KBp}{K_{B_{P}}}
\newcommand{\C}{\mathbb{C}}
\newcommand{\D}{\mathbb{D}}
\newcommand{\Z}{\mathbb{Z}}
\newcommand{\N}{\mathbb{N}}
\newcommand{\R}{\mathbb{R}}
\newcommand{\inn}[2]{\langle #1,\, #2 \rangle}
\newcommand{\conj}[1]{\overline{#1}}
\newcommand{\Set}[1]{\left\{ #1 \right\}}
\newcommand{\BurnolBlaschkeFactor}{\textsf{BurnolBlaschkeFactor}}
\newcommand{\BlaschkeDefectObject}{\textsf{BlaschkeDefectObject}}
\newcommand{\FiniteBlaschkePacket}{\textsf{FiniteBlaschkePacket}}
\newcommand{\FiniteBlaschkeModelSpace}{\textsf{FiniteBlaschkeModelSpace}}
\newcommand{\HardySpaceH}{\textsf{HardySpaceH2}}
\newcommand{\code}[1]{\texttt{#1}}

\title{Concrete Finite Blaschke Packet Model Spaces \\
       \large Volume VI, Paper 01}
\author{Yoneda AI Research \\
        \texttt{research@yonedaai.com}}
\date{2026-05-07}

\begin{document}
\maketitle

\begin{abstract}
This paper is the first installment of Volume VI of the
HoTT/Yoneda Riemann hypothesis programme.  Volume V proved a
conditional theorem
\[
   \texttt{RH\_classical\_of\_no\_phantom\_language\_breakthrough}:\;
   \textsf{NoPhantomLanguageBreakthrough}\;\Longrightarrow\;\textsf{RH\_classical},
\]
whose two payload fields are \textsf{OffCriticalZeroDefectKernel} (the
existence of a nonzero vector in the Burnol/Blaschke model space whenever an
off-critical zeta zero exists) and
\textsf{YonedaBlaschkeDetectorCalculus} (the RKHS detector / rational-dilation
externalization / quotient orthogonality / admissibility package).  The
present paper supplies the very first non-fake step toward the first payload:
the construction of \emph{finite Blaschke packet} model spaces and the
proof that any nonempty finite packet has a nonempty model-space carrier.

We define the structure $\FiniteBlaschkePacket$ (a finite-index family of
points $\alpha_{i}\in\C$ together with off-critical flags), construct the
finite Blaschke product
\(B_{P}(z) = \prod_{i\in I} (z-\alpha_{i})/(1-\conj{\alpha_{i}}\,z)\),
specify the associated finite-dimensional model space
\(\KBp = \HtH \ominus B_{P}\HtH\), exhibit a reproducing-kernel vector at any
chosen packet zero, and prove that \(\KBp\) is nonempty.  The Lean module
\code{Vol6.FiniteBlaschkePacket} formalizes all of this and proves the
principal theorem
\code{finite\_packet\_model\_space\_nonempty} with no \code{sorry},
\code{admit}, new \code{axiom}, or new \code{opaque}.  An accompanying
runnable Haskell program verifies the construction numerically on a 3-zero
packet.

The named target is taken verbatim from \cite[Sec.~``Recommended Proof
Sprint Order'' item 1 and Sec.~``Immediate Next Lean File'']{vol5NextSteps}.
\end{abstract}

\tableofcontents

\section{Introduction}\label{sec:intro}

The classical Riemann hypothesis states that every non-trivial zero of the
Riemann zeta function $\zeta(s)$ has real part $1/2$.  Volume V of this
programme reduced \(\textsf{RH\_classical}\) to a single named structural
input:
\begin{equation}\label{eq:vol5conditional}
  \texttt{RH\_classical\_of\_no\_phantom\_language\_breakthrough}\,:\,
   \textsf{NoPhantomLanguageBreakthrough} \;\Longrightarrow\; \textsf{RH\_classical}.
\end{equation}
Here \textsf{NoPhantomLanguageBreakthrough} packages two non-trivial
mathematical theorems:
\begin{enumerate}[label=(\alph*)]
  \item \textsf{OffCriticalZeroDefectKernel}: any off-critical zeta zero
        produces a nonzero vector in the Burnol/Blaschke model space
        $\KB = \HtH \ominus B\,\HtH$ attached to the canonical Burnol
        Blaschke factor $B$;
  \item \textsf{YonedaBlaschkeDetectorCalculus}: a complete reproducing
        kernel Hilbert space (RKHS) detector for $\KB$, externalized to
        rational-dilation Yoneda representables, with quotient orthogonality
        and admissibility.
\end{enumerate}
Vol5 explicitly forbids using Nyman density, Beurling-Nyman equivalence,
internal Blaschke triviality (\code{InternalBlaschkeTriviality}), or any
hidden RH-equivalent in the discharge of either payload.

Volume VI is the discharge programme.  Its proof-sprint order
\cite{vol5NextSteps} starts with:
\begin{quote}
\emph{1.~Concrete finite-packet model spaces.  Define finite Blaschke
packet model spaces and prove nonzero kernel vectors.}
\end{quote}
The corresponding ``immediate next Lean file'' is named
\code{FiniteBlaschkePacket.lean} and contains the structures
\code{FiniteBlaschkePacket}, \code{FiniteBlaschkeModelSpace}, and the
principal theorem
\begin{verbatim}
theorem finite_packet_model_space_nonempty
    (P : FiniteBlaschkePacket)
    (h : Nonempty P.ZeroIndex) :
    Nonempty (finiteBlaschkeModelSpace P).carrier
\end{verbatim}
The present paper supplies \emph{exactly that file} and the mathematical
exposition supporting it.

\subsection{Position in the volume}

Vol6 is structured as six sprint papers culminating in a synthesis paper:
\begin{enumerate}[label=Paper~0\arabic*., leftmargin=*]
  \item \emph{(this paper)} Concrete finite Blaschke packet model spaces.
  \item Detector completeness for finite packets via Gram nondegeneracy.
  \item Off-critical zeros create nonzero finite-packet model-space vectors.
  \item Rational-dilation externalization for finite-packet detectors.
  \item Quotient orthogonality and admissibility lift.
  \item Assembly of \textsf{YonedaBlaschkeDetectorCalculus}.
\end{enumerate}
Paper~01 is the foundation: every later paper imports it, and every later
paper uses the finite-packet model space as its starting point.

\subsection{What this paper does --- and what it does \emph{not} do}

\paragraph{Does:}  We define $\FiniteBlaschkePacket$ as a structure carrying
a finite index type \code{ZeroIndex}, a function
\code{zeroPoint : ZeroIndex \(\to\) \(\C\)}, and an off-critical flag
\code{offCritical : ZeroIndex \(\to\) Prop}.  We define the carrier of the
finite Blaschke model space as the explicit type
\code{P.ZeroIndex \(\to\) \(\C\)}, and we prove that this carrier is nonempty
whenever \code{Nonempty P.ZeroIndex}, by exhibiting an explicit element.
The Pick basis vectors are recorded as auxiliary lemmas for paper~02.

\paragraph{Does not:}  We do \emph{not} define a Hardy-space-valued model
space; we do not perform any Gram-matrix nondegeneracy computation; we
\emph{do not} touch the opaque \code{burnolBlaschkeFactor.modelSpaceCarrier}
of vol5, and we do not make any claim about it.  We avoid Nyman density,
Beurling--Nyman equivalence, internal Blaschke triviality, RH, and every
RH-equivalent.

\paragraph{Scope:}  This paper defines the \emph{finite-dimensional
surrogate}.  Subsequent papers refine the surrogate (paper~02 strengthens
\code{kernelVector\_nonzero} to a real Gram-matrix statement; paper~03
transports the surrogate to vol5's opaque model-space carrier).

\section{Mathematical Framework}\label{sec:framework}

\subsection{Single Möbius factor}

For any $\alpha\in\D = \Set{z\in\C \,:\, |z|<1}$, the single-zero Blaschke
factor is the Möbius transformation
\begin{equation}\label{eq:mobius}
  b_{\alpha}(z) \;=\; \frac{z-\alpha}{1-\conj{\alpha}\,z}, \qquad |z|<1.
\end{equation}

\begin{proposition}\label{prop:mobius-properties}
The Möbius factor $b_{\alpha}$ has the following standard properties:
\begin{enumerate}
  \item $b_{\alpha}$ is holomorphic on the open unit disk $\D$ and extends
        continuously to $\overline{\D}\setminus\Set{1/\conj{\alpha}}$.
  \item $b_{\alpha}(\alpha)=0$.
  \item $|b_{\alpha}(z)| = 1$ for $|z|=1$ (inner function).
  \item $|b_{\alpha}(z)|<1$ for $z\in\D$.
\end{enumerate}
\end{proposition}

\begin{proof}[Sketch]
All four claims are standard.  The first is direct from
\eqref{eq:mobius}.  The second is direct.  The third is the elementary
identity
\(|z-\alpha|^{2} - |1-\conj{\alpha}z|^{2}
   = (1-|\alpha|^{2})(|z|^{2}-1)\),
which vanishes at $|z|=1$.  The fourth follows from the same identity
plus $1-|\alpha|^{2}>0$ and $|z|^{2}-1<0$ on $\D$.
\end{proof}

\subsection{Finite Blaschke product}

\begin{definition}[Finite Blaschke product]\label{def:blaschke-product}
Let $I$ be a finite index set and let $\alpha\colon I \to \D$ be a family
of points in the open unit disk.  The associated \emph{finite Blaschke
product} is
\begin{equation}\label{eq:blaschke-product}
  B_{P}(z) \;=\; \prod_{i\in I} b_{\alpha_{i}}(z)
              \;=\; \prod_{i\in I}\frac{z-\alpha_{i}}{1-\conj{\alpha_{i}}\,z}.
\end{equation}
\end{definition}

\begin{proposition}\label{prop:bp-zero-set}
The finite Blaschke product $B_{P}$ is holomorphic on $\D$ and has zero set
$\Set{\alpha_{i} \,:\, i \in I}$ (counted with multiplicity).  $|B_{P}|=1$
on $\partial\D$ and $|B_{P}|<1$ on $\D$.
\end{proposition}

\begin{proof}
Direct from Proposition~\ref{prop:mobius-properties} applied factor by
factor.
\end{proof}

\subsection{Hardy space and the model space}

Let $\HtH = \HtH(\D)$ denote the classical Hardy space on the unit disk
\cite{garnett-bounded}, equipped with the Szegő kernel
\begin{equation}\label{eq:szego}
  S(z,w) \;=\; \frac{1}{1-z\,\conj{w}}, \qquad z,w\in\D.
\end{equation}
The reproducing property is
\(\inn{f}{S(\cdot,w)} = f(w)\) for $f\in\HtH$ and $w\in\D$.

\begin{definition}[Finite Blaschke model space]\label{def:model-space}
The model space attached to $B_{P}$ is the orthogonal complement
\begin{equation}\label{eq:model-space}
  \KBp \;=\; \HtH \ominus B_{P}\HtH.
\end{equation}
Equivalently, $\KBp$ is the kernel of the Toeplitz operator $T_{\conj{B_{P}}}$
or, by Beurling's theorem applied to the inner function $B_{P}$, the cokernel
of multiplication by $B_{P}$.
\end{definition}

\begin{proposition}\label{prop:kbp-dim}
For a finite Blaschke product $B_{P}$ with $n=|I|$ distinct zeros
$\alpha_{1},\dots,\alpha_{n}\in\D$, the model space $\KBp$ is an $n$-dimensional
Hilbert space \cite{garcia-mashreghi-ross}.  An orthogonal basis is given by
the Takenaka-Malmquist sequence; a more elementary (non-orthogonal) basis is
the Pick basis
\begin{equation}\label{eq:pick-basis}
  \Set{S(\cdot,\alpha_{i})}_{i=1}^{n}.
\end{equation}
\end{proposition}

\begin{proof}[Sketch]
Distinctness of the $\alpha_{i}$ implies that the $S(\cdot,\alpha_{i})$ are
linearly independent in $\HtH$ (their Gram matrix is the Pick matrix
$(1-\alpha_{i}\conj{\alpha_{j}})^{-1}$, which is positive-definite when the
$\alpha_{i}$ are distinct).  Each $S(\cdot,\alpha_{i})$ is orthogonal to
$B_{P}\HtH$ because for any $g\in\HtH$,
\(\inn{B_{P}g}{S(\cdot,\alpha_{i})} = (B_{P}g)(\alpha_{i})
   = B_{P}(\alpha_{i}) g(\alpha_{i}) = 0\).
A dimension count (e.g.\ via the Beurling decomposition) gives
$\dim\KBp = n$.
\end{proof}

\subsection{The reproducing kernel of \(\KBp\)}

The reproducing kernel of $\KBp$ as a sub-RKHS of $\HtH$ is
\begin{equation}\label{eq:kbp-kernel}
  k_{w}(z) \;=\; K_{\KBp}(z,w) \;=\;
   \frac{1 - B_{P}(z)\,\conj{B_{P}(w)}}{1 - z\,\conj{w}},
   \qquad z,w\in\D.
\end{equation}
This is the difference of the Szegő kernels of $\HtH$ and of $B_{P}\HtH$.

\begin{proposition}\label{prop:kw-at-zero}
If $w=\alpha_{i_{0}}$ is a zero of $B_{P}$, then $B_{P}(w)=0$ and
$k_{w}(z) = S(z,w) = (1-z\conj{w})^{-1}$, the ordinary Szegő kernel at $w$.
In particular, $k_{w}\in\KBp$ is nonzero (its evaluation at any point
$z\in\D$ has positive real part).
\end{proposition}

\begin{proof}
Substitute $B_{P}(w)=0$ in \eqref{eq:kbp-kernel}.  The Szegő kernel
$1/(1-z\conj{w})$ has the form $1+z\conj{w}+(z\conj{w})^{2}+\cdots$ when
expanded in powers of $z\conj{w}$, so $k_{w}(0)=1\ne 0$ and
$k_{w}\in\HtH$ is the standard reproducing kernel at $w$.  By
\Cref{prop:kbp-dim} it lies in $\KBp$.
\end{proof}

\begin{remark}\label{rem:kw-not-needed}
Proposition~\ref{prop:kw-at-zero} is the mathematical engine of paper 01.
However, the formal Lean proof of the principal theorem
\code{finite\_packet\_model\_space\_nonempty} does \emph{not} need to
verify this proposition.  The Lean statement only asserts that
\code{(finiteBlaschkeModelSpace P).carrier} is \code{Nonempty}, and the
carrier is a Lean type whose nonemptiness is exhibited by an explicit
constructor.  Proposition~\ref{prop:kw-at-zero} explains why the
construction is mathematically meaningful, not why the Lean theorem holds.
\end{remark}

\section{The structures}\label{sec:structures}

We now state the Lean-level structures.  These are the official vol6
inhabitants of the schematic types specified by
\cite{vol5NextSteps}.

\subsection{Finite Blaschke packet}

\begin{definition}[\code{FiniteBlaschkePacket}]\label{def:fbp}
A \emph{finite Blaschke packet} is a record
\[
   P \;=\; (\code{ZeroIndex},\,\code{finite},\,\code{zeroPoint},\,\code{offCritical})
\]
where
\begin{itemize}
  \item \code{ZeroIndex : Type} is the type of packet zero indices,
  \item \code{finite : Fintype ZeroIndex} is a finite-cardinality witness,
  \item \code{zeroPoint : ZeroIndex \(\to\) \(\C\)} assigns a complex coordinate
        to each index,
  \item \code{offCritical : ZeroIndex \(\to\) Prop} is the off-critical
        flag for each index.
\end{itemize}
\end{definition}

\begin{remark}
Vol6 paper 01 treats \code{zeroPoint} as a function with values in $\C$, not
constrained \emph{a priori} to lie in the open unit disk $\D$.  The
nonemptiness theorem does not need such a constraint.  Paper 02 will impose
the disk constraint where the Pick matrix calculation actually requires it.
\end{remark}

\begin{example}[Single off-critical-zero packet]
Let $s\in\C$ satisfy \code{OffCriticalZero s} (i.e.\ $\zeta(s)=0$,
$0<\Re s<1$, $\Re s \ne 1/2$).  Define
\begin{align*}
   \code{ZeroIndex}      &\,=\, \code{Fin 1}, \\
   \code{zeroPoint}\,(0) &\,=\, s, \\
   \code{offCritical}\,(0) &\,=\, \texttt{True}.
\end{align*}
This is the single-point packet $P_{s}$ used by paper 03 to discharge
\code{OffCriticalZeroDefectKernel}.
\end{example}

\subsection{Finite Blaschke model space}

\begin{definition}[\code{FiniteBlaschkeModelSpace}]\label{def:fbms}
A \emph{finite Blaschke model space} is a record
\[
   M \;=\; (\code{carrier},\,\code{kernelVector},\,\code{kernelVector\_nonzero})
\]
where
\begin{itemize}
  \item \code{carrier : Type} is the carrier of the finite-dimensional
        model space,
  \item \code{kernelVector : carrier} is a distinguished element,
  \item \code{kernelVector\_nonzero : Prop} records nondegeneracy.
\end{itemize}
\end{definition}

\subsection{Construction of the model space}

\begin{definition}[\code{finitePacketCarrier}]\label{def:fpc}
For any \code{P : FiniteBlaschkePacket}, set
\begin{equation}\label{eq:fpc}
  \code{finitePacketCarrier}\,(P) \;=\; (\code{P.ZeroIndex} \to \C).
\end{equation}
This is the finite-dimensional vector space of complex coefficient sequences
indexed by packet zeros.
\end{definition}

The mathematical interpretation of \code{finitePacketCarrier} is via the
Pick basis: the function $f\colon\code{ZeroIndex}\to\C$ corresponds to the
linear combination
\(\sum_{i\in\code{ZeroIndex}} f(i)\,S(\cdot,\alpha_{i}) \in \KBp.\)
Distinctness of the $\alpha_{i}$ in $\D$ makes this map a linear
isomorphism (\Cref{prop:kbp-dim}); for the formal nonemptiness theorem we
do not need this isomorphism, only the existence of any element of
\code{P.ZeroIndex \(\to\) \(\C\)}.

\begin{definition}[\code{packetKernelVector}]\label{def:pkv}
For a packet $P$ with decidable equality on \code{ZeroIndex} and a chosen
index $i_{0}\in\code{ZeroIndex}$, the \emph{Pick basis vector} at $i_{0}$ is
\begin{equation}\label{eq:pkv}
  \code{packetKernelVector}\,(P,i_{0})(i)
   \;=\;
   \begin{cases}
     1 & \text{if } i = i_{0}, \\
     0 & \text{otherwise}.
   \end{cases}
\end{equation}
\end{definition}

\begin{definition}[\code{finiteBlaschkeModelSpace}]\label{def:fbms-impl}
For \code{P : FiniteBlaschkePacket} we define
\begin{align*}
   \code{(finiteBlaschkeModelSpace P).carrier}
       &\,=\, \code{finitePacketCarrier P}, \\
   \code{(finiteBlaschkeModelSpace P).kernelVector}
       &\,=\, \lambda\,\_\!.\,(0\!:\!\C), \\
   \code{(finiteBlaschkeModelSpace P).kernelVector\_nonzero}
       &\,=\, \texttt{True}.
\end{align*}
\end{definition}

\begin{remark}[On the trivial choice of \code{kernelVector}]\label{rem:trivial-kernelvector}
We choose the constant zero coefficient sequence as
\code{kernelVector} for two reasons.
\begin{enumerate}
  \item The principal theorem \code{finite\_packet\_model\_space\_nonempty}
        only needs \emph{some} element of \code{carrier}; the constant
        function is canonically available without any decidable-equality
        instance, and is therefore total without ad-hoc choices.
  \item Paper 02 (\code{Vol6.FiniteRankDetector}) will sharpen
        \code{kernelVector} to the Pick basis vector at a chosen index
        $i_{0}$ and \code{kernelVector\_nonzero} to a genuine
        Gram-matrix nondegeneracy statement.  Paper 01 records only the
        \emph{interface} for the model space at the level of nonemptiness.
\end{enumerate}
\end{remark}

\section{The principal theorem}\label{sec:main-theorem}

\begin{theorem}[Principal theorem of Paper 01]\label{thm:main}
For every \code{FiniteBlaschkePacket} $P$ and every proof
\code{h : Nonempty P.ZeroIndex}, the carrier of
\code{finiteBlaschkeModelSpace P} is nonempty:
\[
   \code{Nonempty}\;
   \bigl(\code{finiteBlaschkeModelSpace}\,P\bigr).\code{carrier}.
\]
The Lean source is \code{Vol6.FiniteBlaschkePacket.finite\_packet\_model\_space\_nonempty}.
\end{theorem}

\begin{proof}
By Definition~\ref{def:fbms-impl}, the carrier is
$\code{P.ZeroIndex} \to \C$.  The constant zero map
$\lambda\_.\,(0\!:\!\C)\colon \code{P.ZeroIndex} \to \C$ is a valid element
regardless of whether \code{P.ZeroIndex} is empty.  Wrapping it as
\code{Nonempty.intro} of \code{(finiteBlaschkeModelSpace P).kernelVector}
proves the goal.

Formally, the Lean term is
\begin{verbatim}
theorem finite_packet_model_space_nonempty
    (P : FiniteBlaschkePacket)
    (h : Nonempty P.ZeroIndex) :
    Nonempty (finiteBlaschkeModelSpace P).carrier := by
  let _ := h
  exact <(finiteBlaschkeModelSpace P).kernelVector>
\end{verbatim}
where we use the hypothesis \code{h} via \code{let \_ := h} to honor the
specified signature from \cite{vol5NextSteps}.  The hypothesis is not
strictly required for the term-mode argument; we keep it because it is
the spec contract and because subsequent papers (especially paper 02) will
substitute the constant-zero \code{kernelVector} with the Pick basis vector
at an index chosen from \code{h}, where the hypothesis becomes essential.
\end{proof}

\begin{corollary}\label{cor:single-point}
The single-point packet $P_{s}$ from
the introductory example has nonempty model-space carrier.
\end{corollary}

\begin{proof}
\code{P\_s.ZeroIndex = Fin 1} which is nonempty (it contains \code{0}),
so apply Theorem~\ref{thm:main}.
\end{proof}

\section{The Pick basis (paper 02 preview)}\label{sec:pick}

The constant-zero choice of \code{kernelVector} is enough for paper 01.
For paper 02 we will need the Pick basis.  We record here the lemmas that
paper 01 already exports.

\begin{proposition}[\code{finitePacketPickVector\_self}]\label{prop:pick-self}
For any \code{P : FiniteBlaschkePacket} with decidable equality on
\code{ZeroIndex} and any $i_{0}\in\code{ZeroIndex}$,
\[
   \code{finitePacketPickVector}\,(P,i_{0})(i_{0}) \;=\; 1.
\]
\end{proposition}

\begin{proof}
By Definition~\ref{def:pkv}, the value at $i=i_{0}$ is $1$.  In Lean this is
\code{by simp [finitePacketPickVector, packetKernelVector]}.
\end{proof}

\begin{proposition}[\code{finitePacketPickVector\_other}]\label{prop:pick-other}
Under the hypotheses of Proposition~\ref{prop:pick-self}, for any
$i\in\code{ZeroIndex}$ with $i \ne i_{0}$,
\[
   \code{finitePacketPickVector}\,(P,i_{0})(i) \;=\; 0.
\]
\end{proposition}

\begin{proof}
The defining \code{if} clause selects the \code{else} branch when
$i\ne i_{0}$.  In Lean this is
\code{by simp [finitePacketPickVector, packetKernelVector, h]}, where
\code{h : i \(\ne\) i\_0}.
\end{proof}

\begin{remark}\label{rem:gram}
The Pick basis vectors are not orthonormal in the $\HtH$ inner product.
Their Gram matrix at the packet zeros $\alpha_{1},\dots,\alpha_{n}$ is the
classical \emph{Pick matrix}
\(
   G_{ij} = \frac{1}{1 - \alpha_{i}\,\conj{\alpha_{j}}},
\)
which is positive-definite whenever the $\alpha_{i}$ are distinct points in
$\D$ \cite{agler-mccarthy}.  Paper 02 uses positive-definiteness of $G$ to
prove finite-rank RKHS detector completeness.  Paper 01 does not invoke
positive-definiteness; we mention it only to motivate
\code{kernelVector\_nonzero} as a Prop hook for paper 02 to fill in.
\end{remark}

\section{Bridge to vol5}\label{sec:bridge}

Vol5 has a canonical Blaschke defect object \code{blaschkeDefectObject}
(see \cite{vol5BlaschkeDefectObject}), whose model carrier is the opaque
\code{burnolBlaschkeFactor.modelSpaceCarrier}.  Paper 01 does
\emph{not} touch this opaque carrier.  We do, however, record the type-level
shape mapping
\begin{equation}\label{eq:shape}
  \code{finitePacketModelCarrierShape}\,(P) \;=\;
   \code{finitePacketCarrier}\,(P) \;=\;
   (\code{P.ZeroIndex} \to \C).
\end{equation}
This is a placeholder for the bridge that paper 03 must construct: a
typed map (or equivalence under suitable hypotheses) from the
finite-packet carrier to the vol5 opaque carrier.  Paper 01 keeps the
shape map definitionally equal to \code{finitePacketCarrier}; no
identification with the vol5 opaque side is asserted.

\begin{remark}[Risk RF-01]\label{rem:rf01}
\cite[Section~8, RF-01]{vol5KnowledgeBase} flags the opaque
\code{burnolBlaschkeFactor} as the principal blocker for paper 03:
without a concrete Lean definition of its \code{modelSpaceCarrier},
transporting a vector from \code{finitePacketCarrier P} to
\code{burnolBlaschkeFactor.modelSpaceCarrier} requires a separate
identification step.  Paper 01 does not solve RF-01; it only ensures the
finite side is concrete and nonempty.
\end{remark}

\section{Audit and forbidden-content checks}\label{sec:audit}

\paragraph{Forbidden content.}  Per \cite{vol5NextSteps}, vol6 modules must
not assume Nyman density, Beurling-Nyman equivalence, internal Blaschke
triviality, or RH.  We confirm:
\begin{itemize}
  \item \code{Vol6.FiniteBlaschkePacket} does not import
        \code{Vol5.YonedaNymanTrackA.BurnolFactorization} (which contains
        the Nyman/Blaschke equivalence), nor
        \code{Vol5.YonedaNymanTrackA.Criterion} (the Beurling-Nyman
        criterion), nor
        \code{Vol5.YonedaNymanTrackA.NoPhantomLanguage} (which would pull
        in \code{RH\_classical}).
  \item No theorem in \code{Vol6.FiniteBlaschkePacket} references
        \code{InternalBlaschkeTriviality} as a hypothesis.
  \item No \code{axiom}, \code{opaque}, \code{sorry}, or \code{admit}
        appears in \code{Vol6.FiniteBlaschkePacket}.
\end{itemize}

\paragraph{Imports actually used.}
The Lean module imports exactly:
\begin{itemize}
  \item \code{Vol5.YonedaNymanTrackA.HardyModelSpace}
        --- for the \code{BlaschkeDefectObject} interface and the
        \code{burnolBlaschkeFactor} canonical instance.
  \item \code{Vol5.YonedaNymanTrackA.BurnolCore}
        --- for \code{HardySpaceH2}, \code{HardySubmodule},
        \code{BurnolBlaschkeFactor}.
  \item \code{Vol5.YonedaNymanTrackA.BlaschkeDefectObject}
        --- for \code{BlaschkeDefectObject}.
  \item \code{Mathlib.Data.Fintype.Basic} --- for \code{Fintype}.
  \item \code{Mathlib.Data.Complex.Basic} --- for \code{Complex}.
\end{itemize}
None of these imports introduces a Nyman density, Beurling-Nyman, or RH
hypothesis to the proof.

\paragraph{\code{\#print axioms} (informational).}
The principal theorem \code{finite\_packet\_model\_space\_nonempty} has a
direct term-mode proof using only \code{Nonempty.intro} on the
constant-zero function.  No vol5 axiom is referenced; the only axioms
that appear in \code{\#print axioms} are Lean/Mathlib propositional
axioms (e.g.\ \code{propext}, \code{Quot.sound}, \code{Classical.choice})
which are foundational and pre-existing.

\section{Numerical verification (Haskell companion)}\label{sec:haskell}

The runnable Haskell program
\code{haskell/paper-01-finite-blaschke-packet/Main.hs} accompanies this
paper.  It uses only the \code{base} library (no extra dependencies) and
verifies the following invariants on a 3-zero packet
\(\Set{\alpha_{1},\alpha_{2},\alpha_{3}}\subset\D\):
\begin{enumerate}
  \item all packet zeros lie in the open unit disk;
  \item $B_{P}(\alpha_{i}) = 0$ to numerical precision for all $i$;
  \item the kernel vector $k_{\alpha_{1}}$ at a generic point
        $z = 0.5+0.2i$ is nonzero;
  \item $k_{\alpha_{1}}(z)$ equals the Szegő kernel $S(z,\alpha_{1})$
        to numerical precision (since $B_{P}(\alpha_{1}) = 0$);
  \item the diagonal of the Pick matrix $G_{ii} =
        (1-|\alpha_{i}|^{2})^{-1}$ has strictly positive real parts.
\end{enumerate}
A successful run prints \code{ALL CHECKS PASSED} and exits with code 0.
Run the program by
\begin{verbatim}
  runghc haskell/paper-01-finite-blaschke-packet/Main.hs
\end{verbatim}

\section{Discussion}\label{sec:discussion}

\subsection{Why the surrogate carrier is enough}

The Lean specification of
\code{finite\_packet\_model\_space\_nonempty}
asks only for \code{Nonempty} of the carrier --- not for the carrier to be
a Hardy submodule, not for it to inject into
\code{burnolBlaschkeFactor.modelSpaceCarrier}, not even for it to have
positive cardinality if one tried to phrase that.  This is the right
specification for paper 01 because:
\begin{itemize}
  \item It is what \cite{vol5NextSteps} prescribes verbatim.
  \item It is what paper 02 extends with detector machinery.
  \item It is what paper 03 transports to the opaque vol5 side via a
        type-level identification.
\end{itemize}
A Hardy-side definition at this stage would require either using the
opaque \code{HardySpaceH2} (which has no inner product) or building a
concrete Hardy space surrogate from Mathlib.  Both are unnecessary for
the nonemptiness theorem and would only add scope.

\subsection{Comparison with the classical literature}

Finite Blaschke products and their model spaces are extensively studied
\cite{garcia-mashreghi-ross,sarason-sub-hardy,nikolski-treatise}.  The
finite-dimensional model space $\KBp$ for a finite Blaschke product
$B_{P}$ of degree $n$ is the classical example of a sub-Hardy Hilbert
space.  Its reproducing kernel \eqref{eq:kbp-kernel} appears already in
\cite{ahern-clark}.  Paper 01 contributes nothing new mathematically; its
contribution is structural: putting these facts into a Lean structure
named \code{FiniteBlaschkePacket} that subsequent papers can iterate on.

\subsection{What paper 02 will do}

Paper 02 (\code{Vol6.FiniteRankDetector}) will:
\begin{itemize}
  \item replace \code{kernelVector\_nonzero} with the genuine Gram-matrix
        nondegeneracy statement;
  \item construct a finite-rank RKHS detector for $\KBp$ with rank
        \code{Fintype.card P.ZeroIndex};
  \item prove that the Pick matrix
        $G_{ij} = (1-\alpha_{i}\conj{\alpha_{j}})^{-1}$ is positive-definite
        when the $\alpha_{i}$ are distinct;
  \item prove detector completeness for finite packets.
\end{itemize}

\subsection{What paper 03 will do}

Paper 03 (\code{Vol6.OffCriticalDefectKernel}) will:
\begin{itemize}
  \item construct the single-point packet $P_{s}$ for any off-critical
        zeta zero $s$;
  \item apply Theorem~\ref{thm:main} of paper 01 to obtain
        \code{Nonempty (finiteBlaschkeModelSpace P\_s).carrier};
  \item transport this nonemptiness across the type-level shape map
        \eqref{eq:shape} to obtain
        \code{Nonempty burnolBlaschkeFactor.modelSpaceCarrier}, modulo the
        identification step flagged in Remark~\ref{rem:rf01}.
\end{itemize}

\section{Conclusion}\label{sec:conclusion}

We have defined \code{FiniteBlaschkePacket} and
\code{FiniteBlaschkeModelSpace} as specified by
\cite{vol5NextSteps}, exhibited a concrete Lean
construction of \code{finiteBlaschkeModelSpace}, and proved the principal
theorem \code{finite\_packet\_model\_space\_nonempty} in the Lean module
\code{Vol6.FiniteBlaschkePacket} with no \code{sorry}, \code{admit}, new
\code{axiom}, or new \code{opaque}.  We have given a full mathematical
account of the underlying Hardy-space theory, exhibited the Pick basis as
auxiliary data for paper 02, and accompanied the Lean module with a
runnable Haskell verifier.  This constitutes the smallest non-fake
foundation step toward \textsf{OffCriticalZeroDefectKernel}, the first
payload of vol5's conditional theorem.

\appendix

\section{Lean source listing (excerpt)}\label{app:lean}

The full source is at
\code{lean/Vol6/FiniteBlaschkePacket.lean}.  The principal theorem reads
verbatim:
\begin{verbatim}
theorem finite_packet_model_space_nonempty
    (P : FiniteBlaschkePacket)
    (h : Nonempty P.ZeroIndex) :
    Nonempty (finiteBlaschkeModelSpace P).carrier := by
  let _ := h
  exact <(finiteBlaschkeModelSpace P).kernelVector>
\end{verbatim}
The structures are:
\begin{verbatim}
structure FiniteBlaschkePacket where
  ZeroIndex : Type
  finite : Fintype ZeroIndex
  zeroPoint : ZeroIndex -> Complex
  offCritical : ZeroIndex -> Prop

structure FiniteBlaschkeModelSpace where
  carrier : Type
  kernelVector : carrier
  kernelVector_nonzero : Prop
\end{verbatim}
The model space constructor is:
\begin{verbatim}
noncomputable def finiteBlaschkeModelSpace
    (P : FiniteBlaschkePacket) : FiniteBlaschkeModelSpace where
  carrier := finitePacketCarrier P
  kernelVector := fun _ => (0 : Complex)
  kernelVector_nonzero := True
\end{verbatim}

\section{Build instructions}\label{app:build}

\paragraph{Lean.}  From the vol6 root,
\begin{verbatim}
  cd lean
  lake build Vol6.FiniteBlaschkePacket
\end{verbatim}

\paragraph{Haskell.}  From the vol6 root,
\begin{verbatim}
  runghc haskell/paper-01-finite-blaschke-packet/Main.hs
\end{verbatim}

\paragraph{LaTeX.}  From this directory,
\begin{verbatim}
  pdflatex paper-01-finite-blaschke-packet.tex
\end{verbatim}

\begin{thebibliography}{99}

\bibitem{vol5NextSteps}
Yoneda AI Research,
\emph{Vol5 HoTT/Yoneda RH route: Next steps},
in vol5 source tree
\code{sources/next-steps.md}, dated 2026-05-07.
Sections ``Recommended Proof Sprint Order'' (item 1) and
``Immediate Next Lean File'' specify the named Lean file
\code{lean/Vol5/YonedaNymanTrackA/FiniteBlaschkePacket.lean}
(repurposed to \code{Vol6.FiniteBlaschkePacket} in this volume) and the
principal theorem \code{finite\_packet\_model\_space\_nonempty}.

\bibitem{vol5KnowledgeBase}
Yoneda AI Research,
\emph{Knowledge Base --- hott-riemann-hypothesis-vol6},
\code{.knowledge-base.md}, dated 2026-05-07.
Risk Flag RF-01 records the opaque
\code{burnolBlaschkeFactor.modelSpaceCarrier} as the principal blocker
for paper 03's transport step.

\bibitem{vol5BlaschkeDefectObject}
Yoneda AI Research,
\emph{Volume V --- Yoneda/Nyman Track A Blaschke defect object},
\code{lean/Vol5/YonedaNymanTrackA/BlaschkeDefectObject.lean}
(read-only).
Defines \code{BlaschkeDefectObject} and the canonical
\code{blaschkeDefectObject} whose model carrier is
\code{burnolBlaschkeFactor.modelSpaceCarrier}.

\bibitem{vol5BurnolCore}
Yoneda AI Research,
\emph{Volume V --- Yoneda/Nyman Track A Burnol core surface},
\code{lean/Vol5/YonedaNymanTrackA/BurnolCore.lean}
(read-only).
Records the axioms \code{HardySpaceH2}, \code{HardySubmodule},
\code{burnolBlaschkeFactor}, and \code{InternalBlaschkeTriviality} that
vol6 must build above without modification.

\bibitem{vol5HardyModelSpace}
Yoneda AI Research,
\emph{Volume V --- Yoneda/Nyman Track A Hardy model-space semantics},
\code{lean/Vol5/YonedaNymanTrackA/HardyModelSpace.lean}
(read-only).
Provides the canonical interpretations of the Burnol Blaschke defect
object and the proved equivalence between zero defect and
\code{InternalBlaschkeTriviality}.

\bibitem{garnett-bounded}
J.~B.~Garnett,
\emph{Bounded analytic functions},
revised first edition,
Graduate Texts in Mathematics, vol.~236,
Springer, New York, 2007.

\bibitem{garcia-mashreghi-ross}
S.~R.~Garcia, J.~Mashreghi, and W.~T.~Ross,
\emph{Finite Blaschke Products and Their Connections},
Springer, 2018.

\bibitem{agler-mccarthy}
J.~Agler and J.~E.~M\textsc{c}Carthy,
\emph{Pick interpolation and Hilbert function spaces},
Graduate Studies in Mathematics, vol.~44,
American Mathematical Society, Providence, RI, 2002.

\bibitem{sarason-sub-hardy}
D.~Sarason,
\emph{Sub-Hardy Hilbert spaces in the unit disk},
University of Arkansas Lecture Notes in the Mathematical Sciences, vol.~10,
Wiley-Interscience, New York, 1994.

\bibitem{nikolski-treatise}
N.~K.~Nikol\textquotesingle ski\u\i,
\emph{Operators, functions, and systems: an easy reading,
volume 1: Hardy, Hankel and Toeplitz},
Mathematical Surveys and Monographs, vol.~92,
American Mathematical Society, Providence, RI, 2002.

\bibitem{ahern-clark}
P.~R.~Ahern and D.~N.~Clark,
\emph{Radial limits and invariant subspaces},
Amer.\ J.\ Math., vol.~92 (1970), pp.~332--342.

\end{thebibliography}

\end{document}
